% Generated draft: requires mathematical and source review.
\documentclass[11pt]{article}
\usepackage[margin=25mm]{geometry}
\usepackage{fontspec}
\setmainfont{Latin Modern Roman}
\usepackage{amsmath,amssymb,mathtools,unicode-math}
\setmathfont{Latin Modern Math}
\usepackage{enumitem,needspace,xurl,hyperref}
\setlength{\parindent}{0pt}
\setlength{\parskip}{5pt}
\setlength{\emergencystretch}{3em}
\newcommand{\N}{\mathbb{N}}
\newcommand{\Z}{\mathbb{Z}}
\newcommand{\Q}{\mathbb{Q}}
\newcommand{\R}{\mathbb{R}}
\newcommand{\C}{\mathbb{C}}
\newcommand{\F}{\mathbb{F}}
\newcommand{\A}{\mathbb{A}}
\newcommand{\PP}{\mathbb{P}}
\newcommand{\E}{\mathbb{E}}
\newcommand{\eps}{\varepsilon}
\newcommand{\dd}{\,\mathrm{d}}
\begin{document}
\section*{UnsolvedMath ID: 10400168}
\textbf{Source problem code:} catalogue entry
\subsection{Definitions and mathematical statement}
A fusion-rule algebra has a distinguished finite label basis with products described by nonnegative integer multiplicities \([i][j]=\sum_kN_{ij}^k[k]\), together with the required unit and duality data. Compatible \(6j\)-symbols encode associativity data satisfying coherence relations used in a TQFT construction.

Y. Kawahigashi assumes that such data give a TQFT \((V_I,Z_I)\), and that a fusion-rule subalgebra, with the \(6j\)-symbols restricted to it, gives another TQFT \((V_J,Z_J)\). The question is to determine the relation between these two theories. In a based-subalgebra interpretation, the subset \(J\) of labels must have the closure and duality properties required by the stated construction; an arbitrary linear subalgebra is not silently assumed to suffice.

The comparison may concern the surface state spaces and cobordism maps, as well as the closed-manifold invariants, but the source does not prescribe a particular kind of relation. In particular, it does not assert that the smaller theory is a subspace, quotient, direct summand, or naturally embedded subfunctor of the larger one. Such a comparison requires proof compatible with the actual TQFT construction.

The premise is that both theories have been constructed from the given and restricted data. The excerpt does not identify that construction, coefficient field, spherical/ribbon/normalization assumptions, or exact subalgebra convention. It must not be read as a theorem that every fusion-rule algebra or every restriction automatically produces a TQFT.
\subsection{Short English statement}
Given a TQFT and another obtained by restricting its fusion and \(6j\)-data to a fusion-rule subalgebra, determine how the theories are related. The construction and desired form of comparison are unspecified.
\subsection{Sources}
\begin{itemize}
\item Local UnsolvedMath registry, permanent ID 10400168, source code catalogue entry; statement attribution Y. Kawahigashi; proposer field and proposal year absent; updated 2026-08-06.
\item Ohtsuki, \emph{Problems on invariants of knots and 3-manifolds} (2002), Problem 9.4, PDF page 138: \url{https://msp.org/gtm/2002/04/gtm-2002-04-024s.pdf}. Registry access date 2026-07-29; PDF font-structure extraction. Provenance distinguishes 231 numbered bold headings. Present primary-literature status and redistribution terms require review.
\item The supplied summary identifies Kawahigashi's question but provides no exact additional reference or verified comparison theorem.
\end{itemize}
\end{document}
